\documentclass[12pt]{article}

\usepackage[margin=1in]{geometry}
\usepackage{enumitem}
\usepackage{amsmath}
\usepackage{amssymb}
\usepackage{array}
\usepackage{titlesec}

\title{Lecture 7 Notes: \\ Aristotle's \textit{Posterior Analytics}, Book II}
\author{}
\date{}

\begin{document}

\maketitle

\section*{Central Theme}

Aristotle's \textit{Posterior Analytics}, Book II gives an account of inquiry,
definition, cause, and first principles.

Book I asked:

\[
\text{What makes a syllogism scientific demonstration?}
\]

Book II asks:

\[
\text{What are we looking for when we inquire scientifically?}
\]

\[
\boxed{
\textit{Posterior Analytics}, \text{ Book II}
=
\text{Aristotle's theory of inquiry, definition, and causal explanation}
}
\]

Scientific inquiry is the search for the middle term: the cause that explains
what a thing is or why a fact is so.

\section*{1. From Demonstration to Inquiry}

Book I established what demonstration is: a syllogism productive of scientific
knowledge.

Book II asks what we are seeking when we seek knowledge.

\[
\text{Book I: What is demonstration?}
\]

\[
\text{Book II: What do we seek in demonstration?}
\]

The answer is:

\[
\boxed{
\text{We seek the cause, and the cause is the middle term.}
}
\]

\subsection*{Teaching Point}

Book II turns the theory of demonstration into a theory of inquiry.

\[
\textit{Prior Analytics}: \text{middle term as connector}
\]

\[
\textit{Posterior Analytics I}: \text{middle term as cause}
\]

\[
\textit{Posterior Analytics II}: \text{middle term as object of inquiry}
\]

\section*{2. The Four Kinds of Inquiry}

Aristotle says that there are four kinds of questions corresponding to the
things we can know.

\[
\begin{array}{ll}
1. & \text{Whether a connection is a fact} \\
2. & \text{Why the connection holds} \\
3. & \text{Whether a thing exists} \\
4. & \text{What the thing is}
\end{array}
\]

\subsection*{Questions About a Subject and Attribute}

The first two questions concern whether a subject has an attribute and why it
has that attribute.

\[
\text{Does the moon suffer eclipse?}
\]

\[
\text{Why does the moon suffer eclipse?}
\]

\subsection*{Questions About Existence and Essence}

The second pair concerns whether a thing exists and what it is.

\[
\text{Does God exist?}
\]

\[
\text{What is God?}
\]

\subsection*{Whiteboard}

\[
\boxed{
\text{That it is} \rightarrow \text{Why it is}
}
\]

\[
\boxed{
\text{Whether it is} \rightarrow \text{What it is}
}
\]

\subsection*{Teaching Point}

Inquiry begins either with a fact to be explained or with a thing whose
existence and nature are being sought.

\[
\boxed{
\text{Scientific knowledge answers questions of fact, cause, existence, and essence.}
}
\]

\section*{3. Every Inquiry Seeks the Middle}

Aristotle says that all these questions ultimately concern the middle term.

When we ask:

\[
\text{Does the moon suffer eclipse?}
\]

we are really asking:

\[
\text{Is there a cause producing eclipse of the moon?}
\]

When we ask:

\[
\text{Why does the moon suffer eclipse?}
\]

we are asking:

\[
\text{What is that cause?}
\]

\subsection*{Whiteboard}

\[
\text{Does the moon suffer eclipse?}
=
\text{Is there a middle?}
\]

\[
\text{Why does the moon suffer eclipse?}
=
\text{What is the middle?}
\]

\[
\boxed{
\text{All scientific inquiry is inquiry into the middle term.}
}
\]

\subsection*{Teaching Point}

The middle term is no longer merely a logical bridge. It is the object of
scientific search.

\[
\boxed{
\text{To inquire scientifically is to search for the explanatory middle.}
}
\]

\section*{4. The Middle Term Is the Cause}

In the \textit{Prior Analytics}, the middle term connected two extremes.

In the \textit{Posterior Analytics}, the middle term explains why the extremes
are connected.

\[
\text{Subject} \rightarrow \text{Middle} \rightarrow \text{Predicate}
\]

\[
\text{Thing} \rightarrow \text{Cause} \rightarrow \text{Attribute}
\]

\subsection*{Teaching Point}

A scientific middle term is a cause.

\[
\boxed{
\text{The middle term answers the question ``why?''}
}
\]

\section*{5. Essence and Cause}

Aristotle closely links essence and cause.

For many things, to know what the thing is is also to know why the fact occurs.

\subsection*{Eclipse}

Question:

\[
\text{What is eclipse?}
\]

Answer:

\[
\text{The privation of the moon's light by the interposition of the earth.}
\]

But this also answers:

\[
\text{Why does the moon suffer eclipse?}
\]

Answer:

\[
\text{Because the earth shuts out the light.}
\]

\subsection*{Teaching Point}

For many scientific objects, the definition gives the cause.

\[
\boxed{
\text{For many things, to know what it is is to know why it is.}
}
\]

\section*{6. Definition and Demonstration}

Book II raises a central problem:

\begin{quote}
Can the same thing be known both by definition and by demonstration?
\end{quote}

Definition and demonstration have different roles.

\[
\text{Definition} \rightarrow \text{what something is}
\]

\[
\text{Demonstration} \rightarrow \text{that something belongs}
\]

\subsection*{Example}

\[
\text{Triangle has angles equal to two right angles.}
\]

This can be demonstrated.

But having angles equal to two right angles is not the definition of triangle.

\subsection*{Teaching Point}

Definition and demonstration are related, but they are not identical.

\[
\boxed{
\text{Definition reveals essence; demonstration proves connection.}
}
\]

\section*{7. Not Everything Demonstrable Is Definable}

Aristotle argues that not everything demonstrable can be defined.

A demonstrated conclusion often states that an attribute belongs to a subject.

\[
\text{Every triangle has angles equal to two right angles.}
\]

But the demonstrated attribute is not therefore the essence of the subject.

\subsection*{Teaching Point}

A property may belong necessarily without being the definition of the subject.

\[
\boxed{
\text{Necessary property} \neq \text{essence}
}
\]

\section*{8. Not Everything Definable Is Demonstrable}

Aristotle also argues that not everything definable can be demonstrated.

Definitions often function as starting points of demonstration. If every
definition had to be demonstrated, demonstration would fall into infinite
regress.

\[
\text{definition} \rightarrow \text{starting point}
\]

\[
\text{demonstration} \rightarrow \text{proof from starting points}
\]

\subsection*{Teaching Point}

Some definitions are basic and indemonstrable.

\[
\boxed{
\text{Demonstration depends on definitions it does not itself prove.}
}
\]

\section*{9. Why Essence Cannot Simply Be Demonstrated}

Aristotle considers whether there can be a syllogism of a thing's essence.

The difficulty is that a syllogism proves an attribute of a subject through a
middle term. But essence is not merely an attribute added to a subject. It is
what the subject is.

\subsection*{Problem}

To prove the essence of man, one would need a middle term that already contains
the essence of man.

\[
\text{To prove essence, the middle must already contain essence.}
\]

But then the proof assumes what it is supposed to prove.

\[
\boxed{
\text{The proof begs the question.}
}
\]

\subsection*{Teaching Point}

Essence cannot be demonstrated as though it were an ordinary predicate.

\[
\boxed{
\text{Essence is not simply one more attribute attached to a subject.}
}
\]

\section*{10. Division Does Not Prove Essence}

Aristotle criticizes definition by division.

Division proceeds by asking questions such as:

\[
\text{Is man animal or inanimate?}
\]

\[
\text{Is animal terrestrial or aquatic?}
\]

\[
\text{Is terrestrial animal footed or footless?}
\]

But each step is assumed, not inferred.

\subsection*{Teaching Point}

Division may classify, but it does not demonstrate.

\[
\boxed{
\text{Division is useful for finding definitions, but it is not proof.}
}
\]

\section*{11. Why Division Is Not Demonstration}

In a genuine demonstration, the conclusion follows necessarily from the
premises.

In division, the definer asks for concessions and then assembles the formula.
But the full definition does not follow with syllogistic necessity from the
earlier choices.

\subsection*{Example}

A divider may arrive at:

\[
\text{man} = \text{animal, mortal, footed, biped, wingless}
\]

But if asked why this whole formula is the essence of man, division alone does
not supply a demonstration.

\subsection*{Teaching Point}

Division can organize the search, but it cannot by itself prove that the
definition is complete.

\[
\boxed{
\text{Division needs guidance from essence, not merely from classification.}
}
\]

\section*{12. Hypothetical Definitions and Begging the Question}

Aristotle considers whether essence can be proved hypothetically.

One might say:

\[
\text{Suppose the definable form consists of certain peculiar attributes.}
\]

\[
\text{Suppose these are the only such attributes.}
\]

\[
\therefore \text{this is the definition.}
\]

But Aristotle replies that the definable form has already been assumed in the
premises.

\subsection*{Teaching Point}

A definition cannot be smuggled into the premises and then presented as a
conclusion.

\[
\boxed{
\text{If the proof assumes the definable form, it does not prove it.}
}
\]

\section*{13. Definition Does Not Prove Existence}

A definition does not by itself prove that the thing defined exists.

\subsection*{Example}

We may understand the phrase:

\[
\text{goat-stag}
\]

But understanding the phrase does not prove that goat-stags exist.

Likewise, one might understand a definition of circle without that definition
alone proving the existence of circles.

\subsection*{Teaching Point}

Meaning is not existence.

\[
\boxed{
\text{To know what a name means is not yet to know that the thing exists.}
}
\]

\section*{14. Starting Again: A Positive Account}

After raising difficulties, Aristotle begins again.

He argues that to know a thing's essential nature is connected with knowing the
cause of its existence.

If a thing has a cause distinct from itself, that cause can be exhibited through
demonstration.

\[
\text{Essence is not demonstrated directly.}
\]

But:

\[
\text{Essence can be revealed through a demonstration that gives the cause.}
\]

\subsection*{Teaching Point}

Demonstration does not prove essence directly, but it can exhibit essence
through cause.

\[
\boxed{
\text{Causal demonstration reveals definable nature.}
}
\]

\section*{15. The Eclipse Example Revisited}

Let:

\[
A = \text{eclipse}
\]

\[
B = \text{earth's interposition}
\]

\[
C = \text{moon}
\]

To ask whether the moon is eclipsed is to ask whether the defining condition
exists.

\[
\text{Is } B \text{ present?}
\]

If the earth is interposed between sun and moon, then the moon suffers eclipse.

\subsection*{Definition}

\[
\text{Eclipse}
=
\text{privation of the moon's light by the earth's interposition}
\]

\subsection*{Teaching Point}

The cause that explains the fact also supplies the definition.

\[
\boxed{
\text{The cause of eclipse is contained in the definition of eclipse.}
}
\]

\section*{16. The Thunder Example}

Aristotle gives another example: thunder.

Question:

\[
\text{What is thunder?}
\]

Answer:

\[
\text{The noise of fire being quenched in clouds.}
\]

Question:

\[
\text{Why does it thunder?}
\]

Answer:

\[
\text{Because fire is quenched in the clouds.}
\]

\subsection*{Whiteboard}

\[
\text{Thunder} = \text{noise in clouds}
\]

Better:

\[
\boxed{
\text{Thunder} = \text{noise of fire being quenched in clouds}
}
\]

\subsection*{Teaching Point}

Scientific definition is causal definition.

\[
\boxed{
\text{A real definition gives the reason why the thing is what it is.}
}
\]

\section*{17. Three Kinds of Definition}

Aristotle distinguishes several senses of definition.

\[
\begin{array}{ll}
1. & \text{An indemonstrable statement of essence} \\
2. & \text{A quasi-demonstration of essence in another form} \\
3. & \text{The conclusion of a demonstration giving essence}
\end{array}
\]

\subsection*{Nominal and Causal Definition}

Some definitions only state the meaning of a name.

\[
\text{Nominal definition} = \text{what the name means}
\]

Other definitions reveal the cause of a thing's existence.

\[
\text{Causal definition} = \text{why the thing is what it is}
\]

\subsection*{Teaching Point}

Not all definitions are equally scientific.

\[
\boxed{
\text{The best definition is explanatory.}
}
\]

\section*{18. The Four Causes as Middle Terms}

Aristotle says that we think we know scientifically when we know the cause.

He names four causes:

\[
\begin{array}{ll}
1. & \text{The definable form} \\
2. & \text{The antecedent that necessitates a consequent} \\
3. & \text{The efficient cause} \\
4. & \text{The final cause}
\end{array}
\]

Each cause can function as the middle term of a proof.

\[
\boxed{
\text{The middle term can express different kinds of cause.}
}
\]

\subsection*{Teaching Point}

Scientific explanation requires selecting the right kind of cause.

\[
\boxed{
\text{Different questions require different causal middles.}
}
\]

\section*{19. Formal Cause as Middle}

Aristotle asks why the angle in a semicircle is a right angle.

The cause is that it is half of two right angles.

\[
C = \text{angle in a semicircle}
\]

\[
B = \text{half of two right angles}
\]

\[
A = \text{right angle}
\]

\[
C \rightarrow B \rightarrow A
\]

\subsection*{Teaching Point}

The definable form can function as the explanatory middle.

\[
\boxed{
\text{The formal cause explains what the thing is.}
}
\]

\section*{20. Efficient Cause as Middle}

Aristotle gives the example of the Persian War.

Question:

\[
\text{Why did war come upon the Athenians?}
\]

Answer:

\[
\text{Because they raided Sardis with the Eretrians.}
\]

\subsection*{Structure}

\[
C = \text{Athenians}
\]

\[
B = \text{unprovoked raiding}
\]

\[
A = \text{having war waged against them}
\]

\[
C \rightarrow B \rightarrow A
\]

\subsection*{Teaching Point}

The efficient cause explains the originating source of an event.

\[
\boxed{
\text{Efficient cause answers: What produced this?}
}
\]

\section*{21. Final Cause as Middle}

Aristotle gives the example of walking after supper.

Question:

\[
\text{Why walk after supper?}
\]

Answer:

\[
\text{For the sake of health.}
\]

More specifically:

\[
\text{Walking after supper prevents food from rising.}
\]

\[
\text{Preventing food from rising contributes to health.}
\]

\subsection*{Structure}

\[
C = \text{walking after supper}
\]

\[
B = \text{non-regurgitation of food}
\]

\[
A = \text{health}
\]

\[
C \rightarrow B \rightarrow A
\]

\subsection*{Teaching Point}

The final cause explains the end for the sake of which something is done.

\[
\boxed{
\text{Final cause answers: For the sake of what?}
}
\]

\section*{22. Necessity and Final Causality}

Aristotle notes that the same phenomenon may be explained both by necessity and
by an end.

\subsection*{Example}

Light passes through a lantern:

\[
\text{by necessity: because small particles pass through pores}
\]

\[
\text{for an end: so that we may avoid stumbling}
\]

\subsection*{Teaching Point}

Aristotle allows layered explanations.

\[
\boxed{
\text{Necessity and purpose can coexist in explanation.}
}
\]

\section*{23. Nature, Necessity, and Chance}

Aristotle distinguishes outcomes that happen:

\[
\text{always}
\]

\[
\text{for the most part}
\]

\[
\text{by chance}
\]

Natural processes often occur for an end, but they may also involve necessity.

By chance, however, nothing comes to be for an end in the strict explanatory
sense.

\subsection*{Teaching Point}

Scientific explanation must distinguish regular natural order from accidental
coincidence.

\[
\boxed{
\text{Chance events do not have final causes as chance events.}
}
\]

\section*{24. Time and Causal Inference}

Aristotle discusses cause and effect across time.

If the effect is present, the cause is present.

If the effect is past, the cause is past.

If the effect is future, the cause is future.

\subsection*{Eclipse}

\[
\text{The moon was eclipsed because the earth intervened.}
\]

\[
\text{The moon is being eclipsed because the earth is intervening.}
\]

\[
\text{The moon will be eclipsed because the earth will intervene.}
\]

\subsection*{Teaching Point}

The tense of the cause must match the tense of the explained fact.

\[
\boxed{
\text{Past effect requires past cause; present effect requires present cause; future effect requires future cause.}
}
\]

\section*{25. Circular Natural Processes}

Aristotle notes that some natural processes occur in cycles.

\subsection*{Example}

\[
\text{earth is moistened}
\rightarrow
\text{exhalation rises}
\rightarrow
\text{cloud forms}
\rightarrow
\text{rain falls}
\rightarrow
\text{earth is moistened}
\]

\subsection*{Teaching Point}

In cyclic natural processes, causal explanation can move through reciprocal
stages.

\[
\boxed{
\text{Nature sometimes displays circular processes of generation.}
}
\]

\section*{26. What Happens Always and What Happens For the Most Part}

Some things happen always and in every case.

\[
\text{always}
\]

Other things happen as a general rule.

\[
\text{for the most part}
\]

\subsection*{Example}

\[
\text{Not every man grows a beard, but it is the general rule.}
\]

For what happens for the most part, the middle term must also hold for the most
part.

\subsection*{Teaching Point}

The kind of conclusion depends on the kind of middle.

\[
\boxed{
\text{A for-the-most-part conclusion requires a for-the-most-part middle.}
}
\]

\section*{27. How to Obtain a Definition}

Aristotle gives a method for finding definitions.

We select attributes that:

\begin{enumerate}[itemsep=0.25em]
    \item always belong to the thing,
    \item are wider than the thing individually,
    \item do not extend beyond the thing's genus,
    \item together become coextensive with the thing.
\end{enumerate}

\subsection*{Triad Example}

A triad is:

\[
\text{number}
\]

\[
\text{odd}
\]

\[
\text{prime}
\]

Each attribute is wider than triad when taken separately, but together they
define triad.

\[
\boxed{
\text{triad} = \text{number} + \text{odd} + \text{prime}
}
\]

\subsection*{Teaching Point}

A definition gathers essential attributes until they are collectively
coextensive with the thing.

\[
\boxed{
\text{Definition = essential synthesis with the same extension as the subject.}
}
\]

\section*{28. Division as a Useful Search Procedure}

Although division is not demonstration, Aristotle gives it a positive role.

Division helps us:

\[
\begin{array}{ll}
1. & \text{admit only elements of the definable form} \\
2. & \text{arrange them in the right order} \\
3. & \text{omit none of them}
\end{array}
\]

\subsection*{Teaching Point}

Division is not proof, but it is a disciplined way to search for definitions.

\[
\boxed{
\text{Division helps discover essence, though it does not demonstrate it.}
}
\]

\section*{29. The Right Order of Definition}

The order of terms in a definition matters.

It is not the same to say:

\[
\text{animal, tame, biped}
\]

as to say:

\[
\text{biped, animal, tame}
\]

Definition should move from genus to successive differentiae in proper order.

\[
\text{genus} \rightarrow \text{differentia} \rightarrow \text{lower differentia}
\]

\subsection*{Teaching Point}

A definition is not a heap of attributes. It is an ordered account of essence.

\[
\boxed{
\text{Definition requires structure, not mere accumulation.}
}
\]

\section*{30. The Three Rules for Definition by Division}

When establishing a definition by division, Aristotle says we should keep three
things in view.

\[
\begin{array}{ll}
1. & \text{Admit only elements in the definable form} \\
2. & \text{Arrange them in the right order} \\
3. & \text{Omit no essential elements}
\end{array}
\]

\subsection*{Teaching Point}

A good definition is essential, ordered, and complete.

\[
\boxed{
\text{A definition must include all and only the essential parts in the right order.}
}
\]

\section*{31. From Particular Cases to Definition}

Aristotle also recommends beginning from similar individuals and asking what
they have in common.

\subsection*{Example: Pride}

If we inquire into pride, we examine proud people such as:

\[
\text{Alcibiades}
\]

\[
\text{Achilles}
\]

\[
\text{Ajax}
\]

We ask what they have in common as proud.

Suppose we find:

\[
\text{intolerance of insult}
\]

Then we compare other cases, such as Lysander or Socrates, and ask whether
there is one common formula.

\subsection*{Teaching Point}

Definition is discovered by collecting common essential features.

\[
\boxed{
\text{Particulars} \rightarrow \text{common feature} \rightarrow \text{definition}
}
\]

\section*{32. Avoiding Equivocation in Definition}

If inquiry produces not one formula but several, then the term may not name one
thing. It may be equivocal.

\subsection*{Teaching Point}

Definition requires unity. If there is no single formula, there may be no single
thing being defined.

\[
\boxed{
\text{One definition requires one essence.}
}
\]

\section*{33. No Metaphors in Definition}

Aristotle warns against metaphorical expressions in definition.

A definition should be clear, literal, and precise.

\[
\text{metaphor} \neq \text{definition}
\]

\subsection*{Teaching Point}

Definitions must be perspicuous because they serve as starting points for
knowledge.

\[
\boxed{
\text{Scientific definition requires clarity, not ornament.}
}
\]

\section*{34. Selecting Connections for Demonstration}

To formulate what we wish to prove, Aristotle says that we should select
analyses and divisions.

We begin with the common genus.

\[
\text{animal}
\]

Then we ask what belongs to every animal.

Next we move to proximate subgenera.

\[
\text{bird}, \quad \text{horse}, \quad \text{human}
\]

Then we ask what belongs essentially to these.

\subsection*{Whiteboard}

\[
\text{genus}
\rightarrow
\text{proximate subgenus}
\rightarrow
\text{species}
\rightarrow
\text{proper attributes}
\]

\subsection*{Teaching Point}

Scientific inquiry proceeds through ordered classification of genera and
species.

\[
\boxed{
\text{To prove well, first classify well.}
}
\]

\section*{35. Common Characters Beyond Traditional Names}

Aristotle warns that we should not confine ourselves to traditional class names.

We must collect any common character we observe.

\subsection*{Example}

Horned animals may share:

\[
\text{possession of a third stomach}
\]

\[
\text{one row of teeth}
\]

The question then becomes:

\[
\text{To what species does the possession of horns belong?}
\]

\subsection*{Teaching Point}

Scientific classification may require discovering new explanatory groupings.

\[
\boxed{
\text{Science follows real common characters, not merely inherited names.}
}
\]

\section*{36. Analogy in Scientific Selection}

Sometimes there is no single identical name for a common structure.

Aristotle gives examples such as:

\[
\text{a squid's pounce}
\]

\[
\text{a fish's spine}
\]

\[
\text{an animal's bone}
\]

These may have analogous functions or structures, even if they lack one common
name.

\subsection*{Teaching Point}

Analogy can reveal common explanatory structure where language lacks a single
term.

\[
\boxed{
\text{Scientific comparison may discover unity by analogy.}
}
\]

\section*{37. One Middle for Many Connections}

Sometimes one middle explains several connections.

\subsection*{Examples}

\[
\text{echo}
\]

\[
\text{reflection}
\]

\[
\text{rainbow}
\]

These differ specifically, but they are generically alike because they involve
repercussion.

\[
\boxed{
\text{common middle: repercussion}
}
\]

\subsection*{Teaching Point}

One causal structure may appear in different phenomena.

\[
\boxed{
\text{Scientific explanation often discovers a shared cause across diverse cases.}
}
\]

\section*{38. Subordinate Middles}

Some explanations differ because one middle is subordinate to another.

\subsection*{Example}

Question:

\[
\text{Why does the Nile rise toward the end of the month?}
\]

Answer:

\[
\text{Because the month is more stormy toward its close.}
\]

Question:

\[
\text{Why is the month more stormy toward its close?}
\]

Answer:

\[
\text{Because the moon is waning.}
\]

\subsection*{Teaching Point}

Some causes explain through chains of subordinate middles.

\[
\boxed{
\text{A cause may itself require a deeper cause.}
}
\]

\section*{39. Cause and Effect: Must the Cause Be Present?}

Aristotle asks whether, when the effect is present, the cause must also be
present.

If we are dealing with universal and commensurate causal connections, cause and
effect must reciprocate within the proper domain.

\subsection*{Example}

If the cause of deciduousness is coagulation of sap, then:

\[
\text{deciduousness} \leftrightarrow \text{coagulation of sap}
\]

within the appropriate subject domain.

\subsection*{Teaching Point}

Scientific cause must be proper and commensurate with the effect.

\[
\boxed{
\text{A true scientific cause explains all and only the relevant cases.}
}
\]

\section*{40. Cause Is Prior to Effect}

Aristotle rejects the idea that an effect and its cause can simply explain each
other in the same way.

The earth's interposition causes eclipse. Eclipse does not cause the earth's
interposition.

\[
\text{earth's interposition} \rightarrow \text{eclipse}
\]

not:

\[
\text{eclipse} \rightarrow \text{earth's interposition}
\]

\subsection*{Teaching Point}

Cause is prior to effect, even when cause and effect are convertible in a
scientific domain.

\[
\boxed{
\text{Convertibility does not erase explanatory priority.}
}
\]

\section*{41. Can One Effect Have Many Causes?}

Aristotle allows that the same effect may have different causes in different
kinds of subjects.

\subsection*{Example}

Longevity may have one cause in quadrupeds and another cause in birds.

\[
\text{quadrupeds} \rightarrow \text{lack of bile}
\]

\[
\text{birds} \rightarrow \text{dry constitution}
\]

\subsection*{Teaching Point}

Plurality of causes is possible across kinds, but not within the same essential
kind.

\[
\boxed{
\text{Different species may require different causes for similar effects.}
}
\]

\section*{42. The Proximate Cause Is the True Cause}

When there are several middle terms, the true cause is the middle nearest to the
species in which the property is manifested.

\[
\text{remote cause} \rightarrow \text{general explanation}
\]

\[
\text{proximate cause} \rightarrow \text{proper explanation}
\]

\subsection*{Teaching Point}

The best explanation gives the nearest appropriate middle.

\[
\boxed{
\text{Scientific explanation seeks the proximate cause.}
}
\]

\section*{43. The Origin of First Principles}

The final chapter asks how we come to know the primary immediate premises.

Aristotle rejects two views:

\[
\text{First principles are innate in a fully developed form.}
\]

\[
\text{First principles arise from nothing.}
\]

Instead, we have a natural capacity that begins with sense-perception.

\subsection*{Teaching Point}

First principles are neither fully innate nor learned from higher principles.

\[
\boxed{
\text{They arise from perception through a natural cognitive process.}
}
\]

\section*{44. From Sense Perception to Science}

Aristotle gives a developmental account.

\[
\text{sense perception}
\rightarrow
\text{memory}
\rightarrow
\text{experience}
\rightarrow
\text{universal}
\rightarrow
\text{art and science}
\]

\subsection*{Explanation}

Sense-perception grasps particulars.

Repeated perception may persist as memory.

Many memories of the same thing become experience.

From experience, the universal becomes stabilized in the soul.

From this arise art and science.

\subsection*{Teaching Point}

Science begins in perception but culminates in grasping universals.

\[
\boxed{
\text{The universal emerges from repeated experience of particulars.}
}
\]

\section*{45. The Battle-Rout Image}

Aristotle compares the emergence of knowledge to a rout in battle being stopped.

First one soldier makes a stand, then another, and eventually the original
formation is restored.

Likewise, among many perceptions, one stable universal first stands in the soul,
then another, until knowledge becomes ordered.

\subsection*{Teaching Point}

Knowledge arises when unstable perceptions settle into stable universals.

\[
\boxed{
\text{Learning is the stabilization of the universal within experience.}
}
\]

\section*{46. Induction and First Principles}

Aristotle says that we come to know primary premises by induction.

\[
\text{particulars}
\rightarrow
\text{induction}
\rightarrow
\text{universal}
\rightarrow
\text{first principles}
\]

But induction here is not merely counting cases. It is the process through which
the universal becomes present in the soul through repeated perception and
memory.

\subsection*{Teaching Point}

Induction prepares the mind to grasp first principles.

\[
\boxed{
\text{First principles are reached through induction from experience.}
}
\]

\section*{47. Intuition}

Scientific knowledge depends on demonstration, but demonstration depends on
primary immediate premises.

Therefore, scientific knowledge cannot itself be what grasps the first
principles.

Aristotle concludes that the state that grasps first principles is intuition.

\[
\text{intuition} \rightarrow \text{first principles}
\]

\[
\text{first principles} \rightarrow \text{demonstration}
\]

\[
\text{demonstration} \rightarrow \text{science}
\]

\subsection*{Teaching Point}

Intuition is the originative source of scientific knowledge.

\[
\boxed{
\text{Science begins from principles grasped by intuition.}
}
\]

\section*{48. The Architecture of Book II}

Book II moves through the following arc:

\[
\text{four questions}
\rightarrow
\text{middle term as cause}
\rightarrow
\text{definition vs. demonstration}
\rightarrow
\text{critique of proving essence}
\]

\[
\rightarrow
\text{causal definition}
\rightarrow
\text{four causes}
\rightarrow
\text{methods of definition}
\rightarrow
\text{causal selection}
\]

\[
\rightarrow
\text{proximate cause}
\rightarrow
\text{origin of first principles}
\rightarrow
\text{intuition}
\]

\subsection*{Teaching Point}

Book II completes the \textit{Posterior Analytics} by showing that scientific
inquiry seeks causes, definitions, and first principles.

\[
\boxed{
\text{Inquiry is the search for the cause that makes knowledge possible.}
}
\]

\section*{49. Summary of Main Ideas}

\begin{enumerate}[itemsep=0.5em]
    \item Scientific inquiry asks four questions: whether a fact is so, why it is
    so, whether a thing exists, and what it is.

    \item All inquiry seeks the middle term.

    \item The middle term is the cause.

    \item Essence and cause are closely connected.

    \item Definition and demonstration are related but not identical.

    \item Essence cannot simply be demonstrated as if it were an ordinary
    predicate.

    \item Division does not prove essence, but it helps discover definition.

    \item Definition does not prove existence.

    \item Demonstration can exhibit essence by revealing the cause.

    \item Scientific definition is causal definition.

    \item The four causes can function as middle terms.

    \item Necessity and final causality can coexist in explanation.

    \item Definitions are obtained by gathering essential attributes that become
    collectively coextensive with the subject.

    \item Scientific inquiry classifies genera, species, and proper attributes.

    \item A single middle can sometimes explain many connections.

    \item Scientific causes must be proper, prior, and proximate.

    \item First principles arise from sense-perception, memory, experience, and
    induction.

    \item Intuition grasps the first principles from which science demonstrates.
\end{enumerate}

\section*{Final Framing}

\[
\boxed{
\textit{Posterior Analytics}, \text{ Book II, shows that inquiry is the search for cause.}
}
\]

Book I explains what demonstration must be. Book II explains what demonstration
is looking for: the middle term that reveals the cause, the essence, and
ultimately the first principles of science.

\[
\boxed{
\text{Science demonstrates from principles, but intuition grasps the principles.}
}
\]

\end{document}